% GENERATED FILE. Edit canonical lessons, then run npm run generate:books.
% canonical-source-hash: fnv1a64:aa2baf3dd9189617
% canonical-lessons: GU-C135-tarasyu, GU-C135-bimar, GU-C135-taju, GU-C135-paku, GU-C135-kacu, GU-R135-first-pass-animals-and-the-body, GU-R135-second-pass-town-school-and-work

\chapter{Thirsty, Ill, Fresh, Ripe, Raw}
\label{ch:gu-thirsty-ill-fresh-ripe-raw}
\hlchaptermodality{135}

\begin{chapteropening}
\textbf{By the end of this chapter:} \emph{I can say thirsty, ill, fresh, ripe and raw.}

Complete the last lesson of chapter 135: i can say thirsty, ill, fresh, ripe and raw.
\end{chapteropening}

\section[\texorpdfstring{\gu{તરસ્યું} (tarasyũ)}{તરસ્યું (tarasyũ)}]{\gu{તરસ્યું} (tarasyũ) — thirsty}

\label{lesson:GU-C135-tarasyu}



\begin{quote}
Before the new one: say the Gujarati for sad, then the Gujarati for hungry.
\end{quote}

\subsection*{You'll want to know: \gu{તરસ્યું}}
\textbf{\gu{તરસ્યું}} (\emph{tarasyũ}) — \textquotedblleft{}thirsty\textquotedblright{}.

\begin{grammarlens}[title={What you've built}]
One more word for this chapter.
\end{grammarlens}

\subsection*{Your turn}
\begin{itemize}
\raggedright
  \item \emph{Say it:} \emph{tarasyũ}
  \item \emph{Say it:} \emph{tarasyũ}, once more
  \item \emph{Say it:} say \emph{tarasyũ}
  \item \emph{Recall:} say the Gujarati for sad, then the Gujarati for hungry, then say \emph{tarasyũ} again
\end{itemize}

\begin{tcolorbox}[breakable,colback=teal!4,colframe=teal!35!black,title={Before you move on}]
What does \gu{તરસ્યું} mean? (\textquotedblleft{}Thirsty\textquotedblright{}.) Say it once more.
\end{tcolorbox}
\section[\texorpdfstring{\gu{બીમાર} (bīmār)}{બીમાર (bīmār)}]{\gu{બીમાર} (bīmār) — ill}

\label{lesson:GU-C135-bimar}



\begin{quote}
Before the new one: say the Gujarati for hungry, then the Gujarati for thirsty.

Which \emph{we} includes the person you are talking to, \emph{ame} or \emph{āpṇe}? (\emph{Āpṇe}.)
\end{quote}

\subsection*{You'll want to know: \gu{બીમાર}}
\textbf{\gu{બીમાર}} (\emph{bīmār}) — \textquotedblleft{}ill\textquotedblright{}.

\begin{grammarlens}[title={What you've built}]
One more word for this chapter.
\end{grammarlens}

\subsection*{Your turn}
\begin{itemize}
\raggedright
  \item \emph{Say it:} \emph{bīmār}
  \item \emph{Say it:} \emph{bīmār}, once more
  \item \emph{Say it:} say \emph{bīmār}
  \item \emph{Recall:} say the Gujarati for hungry, then the Gujarati for thirsty, then say \emph{bīmār} again
\end{itemize}

\begin{tcolorbox}[breakable,colback=teal!4,colframe=teal!35!black,title={Before you move on}]
What does \gu{બીમાર} mean? (\textquotedblleft{}Ill\textquotedblright{}.) Say it once more.
\end{tcolorbox}
\section[\texorpdfstring{\gu{તાજું} (tājũ)}{તાજું (tājũ)}]{\gu{તાજું} (tājũ) — fresh}

\label{lesson:GU-C135-taju}



\begin{quote}
Before the new one: say the Gujarati for thirsty, then the Gujarati for ill.
\end{quote}

\subsection*{You'll want to know: \gu{તાજું}}
\textbf{\gu{તાજું}} (\emph{tājũ}) — \textquotedblleft{}fresh\textquotedblright{}.

\begin{grammarlens}[title={What you've built}]
One more word for this chapter.
\end{grammarlens}

\subsection*{Your turn}
\begin{itemize}
\raggedright
  \item \emph{Say it:} \emph{tājũ}
  \item \emph{Say it:} \emph{tājũ}, once more
  \item \emph{Say it:} say \emph{tājũ}
  \item \emph{Recall:} say the Gujarati for thirsty, then the Gujarati for ill, then say \emph{tājũ} again
\end{itemize}

\begin{tcolorbox}[breakable,colback=teal!4,colframe=teal!35!black,title={Before you move on}]
What does \gu{તાજું} mean? (\textquotedblleft{}Fresh\textquotedblright{}.) Say it once more.
\end{tcolorbox}
\section[\texorpdfstring{\gu{પાકું} (pākũ)}{પાકું (pākũ)}]{\gu{પાકું} (pākũ) — ripe}

\label{lesson:GU-C135-paku}



\begin{quote}
Before the new one: say the Gujarati for ill, then the Gujarati for fresh.
\end{quote}

\subsection*{You'll want to know: \gu{પાકું}}
\textbf{\gu{પાકું}} (\emph{pākũ}) — \textquotedblleft{}ripe\textquotedblright{}.

\begin{grammarlens}[title={What you've built}]
One more word for this chapter.
\end{grammarlens}

\subsection*{Your turn}
\begin{itemize}
\raggedright
  \item \emph{Say it:} \emph{pākũ}
  \item \emph{Say it:} \emph{pākũ}, once more
  \item \emph{Say it:} say \emph{pākũ}
  \item \emph{Recall:} say the Gujarati for ill, then the Gujarati for fresh, then say \emph{pākũ} again
\end{itemize}

\begin{tcolorbox}[breakable,colback=teal!4,colframe=teal!35!black,title={Before you move on}]
What does \gu{પાકું} mean? (\textquotedblleft{}Ripe\textquotedblright{}.) Say it once more.
\end{tcolorbox}
\section[\texorpdfstring{\gu{કાચું} (kācũ)}{કાચું (kācũ)}]{\gu{કાચું} (kācũ) — raw, unripe}

\label{lesson:GU-C135-kacu}



\begin{quote}
Before the new one: say the Gujarati for fresh, then the Gujarati for ripe.
\end{quote}

\subsection*{You'll want to know: \gu{કાચું}}
\textbf{\gu{કાચું}} (\emph{kācũ}) — \textquotedblleft{}raw, unripe\textquotedblright{}.

\begin{grammarlens}[title={What you've built}]
One more word for this chapter.
\end{grammarlens}

\subsection*{Your turn}
\begin{itemize}
\raggedright
  \item \emph{Say it:} \emph{kācũ}
  \item \emph{Say it:} \emph{kācũ}, once more
  \item \emph{Say it:} say \emph{kācũ}
  \item \emph{Recall:} say the Gujarati for fresh, then the Gujarati for ripe, then say \emph{kācũ} again
\end{itemize}

\begin{tcolorbox}[breakable,colback=teal!4,colframe=teal!35!black,title={Before you move on}]
What does \gu{કાચું} mean? (\textquotedblleft{}Raw, unripe\textquotedblright{}.) Say it once more.
\end{tcolorbox}
\section[(dialogue)]{Animals and the body — a first pass}

\label{lesson:GU-R135-first-pass-animals-and-the-body}



\begin{quote}
Say the words for ripe and raw.
\end{quote}

\subsection*{Your turn}
Say these aloud:

\begin{itemize}
\raggedright
  \item an office, a factory, a hotel, a library, a town square
  \item a lane, a park, a gate, a pen, a pencil
  \item a class, a teacher, a job, a salary, a tailor
  \item low, wide, thin, wet, dry
\end{itemize}

\begin{tcolorbox}[breakable,colback=teal!4,colframe=teal!35!black,title={Before you move on}]
An office? (\textbf{\gu{કચેરી}}.) A factory? (\textbf{\gu{કારખાનું}}.) A hotel? (\textbf{\gu{હોટેલ}}.) A library? (\textbf{\gu{પુસ્તકાલય}}.) A town square? (\textbf{\gu{ચોક}}.) A lane? (\textbf{\gu{ગલી}}.) A park? (\textbf{\gu{બાગ}}.) A gate? (\textbf{\gu{દરવાજો}}.) A pen? (\textbf{\gu{પેન}}.) A pencil? (\textbf{\gu{પેન્સિલ}}.) A class? (\textbf{\gu{વર્ગ}}.) A teacher? (\textbf{\gu{ગુરુ}}.) A job? (\textbf{\gu{નોકરી}}.) A salary? (\textbf{\gu{પગાર}}.) A tailor? (\textbf{\gu{દરજી}}.) Low? (\textbf{\gu{નીચું}}.) Wide? (\textbf{\gu{પહોળું}}.) Thin? (\textbf{\gu{પાતળું}}.) Wet? (\textbf{\gu{ભીનું}}.) Dry? (\textbf{\gu{સૂકું}}.)
\end{tcolorbox}
\section[(dialogue)]{Town, school and work — a second pass}

\label{lesson:GU-R135-second-pass-town-school-and-work}



\begin{quote}
Say the words for clean and dirty.
\end{quote}

\subsection*{Your turn}
Say these aloud:

\begin{itemize}
\raggedright
  \item clean, dirty, sour, spicy, soft
  \item hard, easy, difficult, clever, lazy
  \item brave, calm, happy, sad, hungry
  \item thirsty, ill, fresh, ripe, raw
\end{itemize}

\begin{tcolorbox}[breakable,colback=teal!4,colframe=teal!35!black,title={Before you move on}]
Clean? (\textbf{\gu{ચોખ્ખું}}.) Dirty? (\textbf{\gu{ગંદું}}.) Sour? (\textbf{\gu{ખાટું}}.) Spicy? (\textbf{\gu{તીખું}}.) Soft? (\textbf{\gu{નરમ}}.) Hard? (\textbf{\gu{કઠણ}}.) Easy? (\textbf{\gu{સહેલું}}.) Difficult? (\textbf{\gu{અઘરું}}.) Clever? (\textbf{\gu{હોશિયાર}}.) Lazy? (\textbf{\gu{આળસુ}}.) Brave? (\textbf{\gu{બહાદુર}}.) Calm? (\textbf{\gu{શાંત}}.) Happy? (\textbf{\gu{ખુશ}}.) Sad? (\textbf{\gu{ઉદાસ}}.) Hungry? (\textbf{\gu{ભૂખ્યું}}.) Thirsty? (\textbf{\gu{તરસ્યું}}.) Ill? (\textbf{\gu{બીમાર}}.) Fresh? (\textbf{\gu{તાજું}}.) Ripe? (\textbf{\gu{પાકું}}.) Raw? (\textbf{\gu{કાચું}}.)
\end{tcolorbox}
